% CPSY 1291 -- Recitation handout 7
% Compile:  latexmk -lualatex handout-07-midterm-review.tex
\documentclass[11pt]{article}
\usepackage{cpsy1291-handout}

\begin{document}

\handouthead{7}{Midterm review}
  {Lectures 1--8, in eight pages and fourteen problems}
  {Week 7 \ $\cdot$\ the midterm is in class, Thu 10/22}

\section*{How to use this}
\addcontentsline{toc}{section}{How to use this}

This is a review, not new material. Everything here has appeared in a lecture,
an assignment, or an earlier recitation.

Work in this order, and do not reverse it. \textbf{First} read Section~1, which
is the whole of Lectures 1--8 compressed into the statements you should be able
to make without notes. \textbf{Then} do the problems in Section~2 with the
summary closed. Reading a summary produces the feeling of understanding;
answering a question produces the fact of it, and the gap between those two is
where midterm marks are lost.

\begin{keybox}
\ask{What the exam can ask.} Every lecture in this course does three things:
introduces an \emph{object} (a distance, a matrix, a rule), states \emph{what
it is for}, and states \emph{when it misleads you}. Exam questions come from all
three, and the third is the one students prepare least and are asked about
most.
\end{keybox}

\section{Lectures 1--8, compressed}

\subsection{L01 --- Representations, spaces, and metrics}

A representation is a pattern that stands for something. Put $n$ stimuli down
the rows and $d$ units across the columns and you have the \textbf{response
matrix} $\boldsymbol{X}$, which is the object the first half of this course
operates on.

The \textbf{dot product} $\boldsymbol{a}\cdot\boldsymbol{b} = \sum_i a_i b_i =
\|\boldsymbol{a}\|\|\boldsymbol{b}\|\cos\theta$ is a measurement of
$\boldsymbol{b}$ along the direction $\boldsymbol{a}$.

Three distances, and why there are three:

\begin{center}
\begin{tabular}{@{}lll@{}}
\toprule
Distance & Sensitive to & Blind to \\
\midrule
Euclidean & magnitude and pattern & nothing \\
cosine & pattern only & overall scale \\
correlation & pattern after centering & scale \emph{and} additive offset \\
\bottomrule
\end{tabular}
\end{center}

Correlation distance is cosine distance on row-centered vectors. It is
\emph{undefined} for a constant row --- a dead unit --- because that row has zero
length after centering.

The \textbf{metric axioms}: (i) $d(a,b) = 0$ iff $a = b$ (minimality);
(ii) $d(a,b) = d(b,a)$ (symmetry); (iii) $d(a,c) \le d(a,b) + d(b,c)$ (the
triangle inequality). Cosine and correlation distance violate the triangle
inequality. Human similarity judgments violate all three: people rate ``North
Korea is like China'' higher than the reverse, and identical twins are not
judged maximally similar to themselves.

\subsection{L02 --- Similarity, MDS, and clustering}

A \textbf{dissimilarity matrix} (RDM) is $n \times n$, symmetric, zero
diagonal. It discards the units and keeps only the geometry, which is why 50,000
voxels and 4,096 network units become directly comparable.

\textbf{MDS} finds coordinates whose distances reproduce a given dissimilarity
matrix. Metric MDS fits the distances; non-metric MDS fits only their
\emph{rank order}, which is the weaker and often more defensible assumption for
behavioral data. Three warnings when reading an MDS map: the axes have no
inherent meaning, the whole map can be rotated or reflected freely, and a
2-D solution may be a poor summary of the matrix that produced it.

\textbf{Shepard's universal law}: perceived similarity falls exponentially with
distance in psychological space, $s = a\,e^{-b d}$. It is a law about
\emph{psychological} space, so applying it to raw network features is not a test
of it.

\textbf{Hierarchical clustering} always returns a tree, whether or not the data
has one. Choosing where to cut chooses a level of description, not a true
number of clusters. The linkage rule (single, complete, average) changes the
answer, so it must be reported.

\textbf{Comparing two RDMs}: Spearman correlation on the upper triangle,
\texttt{np.triu\_indices(n, k=1)}. The pairs are not independent, so the
accompanying $p$-value is not trustworthy.

\subsection{L03 --- PCA and eigendecomposition}

An image is a vector of pixels. High-dimensional spaces are mostly empty and
distances in them lose contrast --- the \textbf{curse of dimensionality} --- but
real data occupies a far smaller subspace than it has room for, which is what
makes dimensionality reduction possible at all.

Center the data, form the covariance $\boldsymbol{C} =
\tfrac{1}{n-1}\tilde{\boldsymbol{X}}^{\!\top}\tilde{\boldsymbol{X}}$, take its
eigenvectors. Each eigenvector is a direction $\boldsymbol{C}$ stretches without
rotating ($\boldsymbol{C}\boldsymbol{v} = \lambda\boldsymbol{v}$); each
eigenvalue says how much variance lies along it. PC1 is the direction of
greatest spread. Components are orthogonal and ordered.

\textbf{Projection} gives coordinates in the new basis;
\textbf{reconstruction} with $K$ components is the best rank-$K$ approximation
in squared error. Eigenfaces are exactly this on face images.

Three cautions: PCA is \textbf{linear}, so a curved manifold is not captured;
it is sensitive to \textbf{scale}, so features in different units should be
standardized first; and \textbf{centering is not optional}, since without it PC1
usually just points at the mean.

\subsection{L04 --- $t$-SNE and UMAP}

$t$-SNE converts distances to neighbor probabilities in high dimensions,
builds a heavy-tailed version in two dimensions, and moves the points to match
the two distributions. \textbf{Perplexity} is roughly the effective number of
neighbors each point tries to keep. The heavy tail exists to solve the
\textbf{crowding problem}: two dimensions do not have enough room for all the
moderate distances that exist in high dimensions.

How to misread it, in three ways that all appear in published figures: cluster
\emph{sizes} mean nothing, distances \emph{between} clusters usually mean
nothing, and apparent clusters can appear in data with none. One picture is
never the result --- vary the seed and the perplexity.

UMAP is faster, tends to preserve more global structure, and shares all the
same warnings.

\subsection{L05 --- The geometry of neural representations}

A stimulus condition is not a point but a \textbf{manifold}: repeated
presentations, viewpoints, and nuisance variation trace out a cloud.

\textbf{Effective dimensionality} summarizes the eigenspectrum in one number ---
for instance the participation ratio $(\sum_i \lambda_i)^2 / \sum_i \lambda_i^2$.
It is \emph{estimated}, not observed, and it depends on the stimulus set and the
number of samples as much as on the system.

\textbf{Mixed selectivity}: units that respond to combinations of variables
raise the dimensionality of the population code and make more categorizations
linearly separable. The cost is abstraction --- high-dimensional codes generalize
less readily. That trade-off is the lecture's core claim.

The \textbf{metric zoo} --- RSA, CCA, CKA, shape metrics --- exists because two
systems have no shared coordinates. The metrics disagree with one another, which
is a live problem, not a detail.

\subsection{L06 --- The model neuron}

$a = g(\boldsymbol{w}\cdot\boldsymbol{x} + b)$. The weights are a direction, so a
neuron \textbf{measures how much of its preferred pattern is present}; the bias
shifts the threshold; $g$ is the nonlinearity, usually ReLU.

Because $\boldsymbol{w}\cdot\boldsymbol{x} =
\|\boldsymbol{w}\|\|\boldsymbol{x}\|\cos\theta$, a strong response has two
possible causes --- an intense stimulus or a well-matched one --- and separating
them requires normalizing the input. Among inputs of fixed length, the best
stimulus for a linear unit is its own weight vector.

A \textbf{tuning curve} is response against a stimulus parameter. A single
linear-threshold unit computes AND, OR, and NOT, but not XOR.

\subsection{L07 --- How one neuron learns}

\begin{center}
\begin{tabular}{@{}llll@{}}
\toprule
Rule & Update & Needs a teacher? & Converges to \\
\midrule
Hebb   & $\Delta\boldsymbol{w} = \eta\,a\,\boldsymbol{x}$ & no & PC1's direction, but $\|\boldsymbol{w}\|$ diverges \\
Oja    & $\Delta\boldsymbol{w} = \eta\,a\,(\boldsymbol{x} - a\boldsymbol{w})$ & no & PC1, unit length \\
Perceptron & $\Delta\boldsymbol{w} = \eta\,(y-\hat{y})\,\boldsymbol{x}$, on errors only & yes & \emph{some} separating boundary \\
Delta  & $\Delta\boldsymbol{w} = \eta\,(y-\hat{y})\,\boldsymbol{x}$, every example & yes & the least-squares solution \\
\bottomrule
\end{tabular}
\end{center}

Hebbian learning finds PC1 because the average weight change is
$\boldsymbol{C}\boldsymbol{w}$, and repeatedly applying $\boldsymbol{C}$ amplifies
the top eigenvector fastest. \textbf{Centering matters}: on uncentered data the
rule chases the mean instead. Oja's rule adds one decay term that normalizes the
weight, and that is the entire difference.

The perceptron and the delta rule share an update but not a behavior: the
perceptron updates only on mistakes and stops at the first boundary that
separates, so \emph{which} boundary you get depends on the order of
presentation. The delta rule updates always, minimizes squared error, and
converges to one answer regardless of order. The perceptron convergence theorem
bounds the number of updates by $(R/\gamma)^2$, where $\gamma$ is the margin ---
so a small margin means many updates, and non-separable data means it never
stops.

\textbf{Rescorla--Wagner} is the delta rule in animal learning: $\Delta V_i =
\eta\,(\lambda - \sum_j V_j)$. Because the error uses the \emph{sum} of all
cues' predictions, a cue already predicted by another cue learns nothing ---
which is \textbf{blocking}, and it is why the rule was a landmark.

\subsection{L08 --- Multilayer networks and backpropagation}

Stacking linear layers gives a linear map, so \textbf{the nonlinearity is not
optional}. XOR is solvable by a 2--2--1 network: the hidden layer re-represents
the four points so that they become linearly separable, which is the whole idea
of a hidden layer in one example.

\textbf{Universal approximation}: one hidden layer of sufficient width can
approximate any continuous function on a bounded domain to any accuracy. It
says nothing about how wide, and nothing about whether gradient descent can
\emph{find} the weights. Depth is about efficiency, not possibility.

\textbf{Backpropagation} is the chain rule applied backwards through the layers,
reusing each layer's stored quantities. It solves credit assignment. It is not
an optimizer --- gradient descent is; backprop only supplies the gradient.

Batch (all data per step), minibatch (a subset), stochastic (one example).
The learning rate is the main knob: too small is slow, too large oscillates or
diverges.

\section{Problems}

Do these closed-book. Answers follow.

\ask{1.} A unit's activity is $(2, 4, 6)$ across three stimuli; another's is
$(1, 2, 3)$. Give the cosine distance and the correlation distance between them,
and explain the pair of answers in one sentence.

\ask{2.} Why is correlation distance undefined for a unit that produces the same
value for every stimulus?

\ask{3.} You are told $d(A,B) = 0.2$, $d(B,C) = 0.2$, $d(A,C) = 0.9$. Which
axiom fails, and name a distance used in this course that permits it.

\ask{4.} Someone reports that an MDS map's horizontal axis is ``animacy''. What
is the methodological objection, and what evidence would answer it?

\ask{5.} You forget to center your data before PCA on a set of images. Describe
PC1.

\ask{6.} A representation gives 90\% of variance in 8 components. The stimulus
set is 40 objects at 30 viewpoints. Why is ``this representation is
8-dimensional'' unsafe?

\ask{7.} Two $t$-SNE runs on identical data, differing only in
\texttt{random\_state}, give visibly different maps. Is either wrong? What
should be reported?

\ask{8.} State the trade-off that mixed selectivity buys and pays for.

\ask{9.} A linear unit has $\boldsymbol{w} = (3, 4)$ and $b = 0$. Give its
response to $\boldsymbol{x} = (1, 0)$, and to the unit-length input that
maximizes its response.

\ask{10.} Hebb's rule is run on data with a large positive mean. What does the
weight converge toward, and what one preprocessing step fixes it?

\ask{11.} Write Oja's rule and say precisely what the extra term relative to
Hebb's rule accomplishes.

\ask{12.} The perceptron is trained twice on the same separable data, differing
only in the order of examples, and gives two different boundaries. Explain, and
say whether the delta rule would do the same.

\ask{13.} In Rescorla--Wagner, cue A is trained to predict a reward for 50
trials. Then A and B are presented together for 50 more, still rewarded. What
happens to $V_B$, and what is this phenomenon called?

\ask{14.} A 2--2--1 network fails to learn XOR from some initializations. Give
one concrete mechanism, and one diagnostic you could run on a failed network.

\subsection*{Answers}

\ask{1.} Cosine distance $0$: the two vectors point in exactly the same
direction, differing only by a factor of two. Correlation distance $0$ as well:
after centering each row, both become proportional to $(-1, 0, 1)$. The pair of
answers says these units carry the same information about the stimuli, differing
only in gain --- which is why neither distance can tell them apart, and why
Euclidean distance would.

\ask{2.} Centering a constant row gives the zero vector, whose norm is zero, and
the correlation is a division by that norm. In code it surfaces as \texttt{nan},
which then propagates silently through the whole matrix.

\ask{3.} The triangle inequality: $0.9 > 0.2 + 0.2$. Cosine distance and
correlation distance both permit this.

\ask{4.} MDS axes carry no inherent meaning --- the solution can be rotated or
reflected freely, so ``the horizontal axis'' is not a property of the data.
Answering it requires independent evidence: correlate the coordinate with an
external animacy rating, or show that a classifier trained on that coordinate
alone predicts animacy on held-out stimuli.

\ask{5.} PC1 points approximately at the mean image, because the uncentered
second-moment matrix is dominated by it. You have spent your first component
describing the average, and the interesting structure moves to PC2 onward.

\ask{6.} The 1,200 images come from only 40 distinct objects; the 30 viewpoints
of one object trace a smooth low-dimensional curve rather than adding
independent directions. The figure is partly a property of how the stimulus set
was built and would move if you used 10 viewpoints with an unchanged system.

\ask{7.} Neither is wrong; both are members of a family generated by the
algorithm's random initialization. Report what is stable across seeds (and
across perplexities), and say how many you ran.

\ask{8.} It buys \textbf{separability}: with units tuned to combinations, more
categorizations become linearly readable, and the effective dimensionality
rises. It pays in \textbf{abstraction}: a high-dimensional code generalizes less
readily to new conditions, because it encodes conjunctions rather than factors.

\ask{9.} $\boldsymbol{w}\cdot(1,0) = 3$. The maximizing unit-length input is
$\boldsymbol{w}/\|\boldsymbol{w}\| = (0.6, 0.8)$, giving a response of
$\|\boldsymbol{w}\| = 5$.

\ask{10.} Toward the mean direction of the data rather than the direction of
greatest \emph{variance} --- Hebbian dynamics follow the second-moment matrix,
which equals the covariance only after centering. Subtract the mean.

\ask{11.} $\Delta\boldsymbol{w} = \eta\,a\,(\boldsymbol{x} - a\boldsymbol{w})$.
The $-\eta a^2 \boldsymbol{w}$ term is a decay proportional to the unit's own
activity: it grows exactly when the weight grows, driving $\|\boldsymbol{w}\|$
toward 1 and preventing the unbounded growth of plain Hebbian learning, without
changing the direction it converges to.

\ask{12.} The perceptron updates only when it makes a mistake, and stops as soon
as no mistakes remain --- so it halts at whichever separating boundary it reaches
first, and which one that is depends on the order the mistakes arrived in. The
delta rule minimizes a squared error with a single minimum and would converge to
the same boundary regardless of order.

\ask{13.} $V_B$ stays near zero. The prediction error is $\lambda - (V_A + V_B)$
and $V_A$ already accounts for the reward, so there is no error left to drive
learning about B. This is \textbf{blocking}. (Remove phase 1 and both cues learn,
sharing the association.)

\ask{14.} Mechanism: a hidden unit's pre-activation is negative for all four
inputs, so its ReLU gradient is zero and it never updates --- a dead unit,
leaving the network effectively 2--1--1, which cannot solve XOR. Diagnostic:
run the four inputs through the trained network and count, for each hidden unit,
how many inputs it responds to with a non-zero value. A unit answering zero for
all four is dead. (Also acceptable: both hidden units converging to the same
function, or a saddle point.)

\section*{The night before}
\addcontentsline{toc}{section}{The night before}

Six things worth being able to state cold, because each has appeared in a
lecture, an assignment, and a recitation:

\begin{enumerate}
  \item the three distances and what each is blind to;
  \item the three metric axioms, and one real violation of each;
  \item what an eigenvector of the covariance is, in a sentence with no algebra;
  \item why Hebb needs centering and what Oja's extra term does;
  \item why the perceptron's answer depends on presentation order and the delta
        rule's does not;
  \item why a stack of linear layers is a linear layer.
\end{enumerate}

\end{document}
